\section{Dashboard}

GIO dient hauptsachlich der Visualisierung von Daten für Mitarbeiter:innen 
damit diese schnell und einfach den Zustand der Infrastruktur erkennen können. 
Die Visualisierung erfolgt über ein Dashboard, 
das die gesammelten Daten in einer übersichtlichen und interaktiven Form darstellt. 
Das Dashboard bietet verschiedene Funktionen, die es den Benutzer:innen ermöglichen, 
die Informationen nach ihren Bedürfnissen zu filtern und zu analysieren.


\subsection{Icons}

Auf der Weltkarte werden die Standorte der einzelnen Standorte durch Icons dargestellt.
\todo{Hier könnte ein Bild der Legende den Icons eingefügt werden.}
Dadurch können die Benutzer:innen schnell erkennen um welche Art von Standort es sich handelt.


\subsection{Gerätestatus}

Da die Geräte an den Standorten unterschiedliche Zustände haben können,
werden diese Zustände durch unterschiedliche Farben der Icons dargestellt.
Diese Farben geben den Benutzer:innen einen schnellen Überblick über den Zustand der Geräte an den Standorten.
Der Status der Geräte kann dabei in drei Kategorien unterteilt werden:
\begin{itemize}
    \item \textbf{Grün:} Das Gerät ist in einem guten Zustand und funktioniert einwandfrei.
    \item \textbf{Gelb:} Das Gerät hat einen Warnstatus, der auf ein mögliches Problem hinweist.
    \item \textbf{Rot:} Das Gerät hat einen kritischen Status, der auf ein ernsthaftes Problem hinweist.
\end{itemize}


\subsection{Detailebenen}
Um die numerischen Informationen auf dem Dashboard mit den Standorten zu verknüpfen, 
wird bei Klick auf ein Icon die detaillierte Ansicht des jeweiligen Standorts geöffnet.
Diese detaillierte Ansicht wird bestimmt durch das zoom level des Dashboards.
Es gibt drei Detailebenen, die jeweils unterschiedliche Informationen anzeigen:
\begin{itemize}
    \item \textbf{Zoom Level 1:} Auf dieser Ebene werden die Standorte auf einer Weltkarte angezeigt. 
    \item \textbf{Zoom Level 2:} Auf dieser Ebene werden die Standorte auf einer Landkarte angezeigt. 
    \item \textbf{Zoom Level 3:} Auf dieser Ebene werden die Standorte auf einer Straßenkarte angezeigt. 
    \todo{Hier muss noch eine Beschreibung der Detailebenen hinzugefügt werden.}
\end{itemize}

